\chapter{Metodología y Planificación}
\label{chap:metodologia}

\lettrine{E}{l} desarrollo de un proyecto de software de la envergadura de un Trabajo de Fin de Grado requiere una organización rigurosa para garantizar el cumplimiento de los plazos y la calidad de los entregables. 

\section{Metodología de Desarrollo}
Debido a la naturaleza investigadora y de desarrollo individual de este proyecto, no se ha adoptado una metodología ágil estricta (como Scrum), que está más orientada a equipos de desarrollo con roles definidos. 

En su lugar, se ha optado por un enfoque \textit{iterativo e incremental}. Este enfoque permite dividir el ciclo de vida del proyecto en ciclos más cortos y manejables (iteraciones), en cada uno de los cuales se añade funcionalidad o se refinan las características existentes.

Las fases iterativas de este proyecto han sido las siguientes:
\begin{enumerate}
    \item \textbf{Análisis e investigación teórica:} Definición de los factores de riesgo de abandono escolar y estudio de la viabilidad de extracción de datos.
    \item \textbf{Diseño de la Arquitectura:} Elección de la pila tecnológica y diseño del ecosistema de contenedores Docker.
    \item \textbf{Modelado de Datos:} Diseño del esquema en estrella y mapeo origen-destino.
    \item \textbf{Desarrollo ETL:} Implementación iterativa de los pipelines (extracción de tablas transaccionales, transformación y carga dimensional).
    \item \textbf{Desarrollo Backend y Frontend:} Exposición de los datos mediante una API RESTful y creación de los \textit{dashboards} visuales, ajustando los umbrales de alerta en base a los resultados observados.
\end{enumerate}

Esta estrategia ha dotado al desarrollo de la flexibilidad necesaria para ajustar las reglas de negocio, como las puntuaciones de riesgo por necesidades educativas especiales o los umbrales socioeconómicos, a medida que se profundizaba en la casuística de los datos reales.

\section{Gestión y Control de Versiones}
Para el correcto seguimiento de la evolución del proyecto, se ha empleado \textbf{Git} como sistema de control de versiones descentralizado. El repositorio central se ha alojado en GitHub, lo que ha facilitado mantener un histórico organizado de los cambios estructurales.

Siguiendo las mejores prácticas de la industria, el repositorio se estructuró monorepo separando claramente los dominios funcionales en directorios independientes (\textit{api}, \textit{dwh}, \textit{hop}, \textit{frontend}), lo que permitió aplicar la metodología iterativa sobre componentes aislados sin interferir con el resto del ecosistema.
